%% DeReFusion Ablation Study Report (English, IEEEtran conference, two-column)
%% Compile: pdflatex main_ieee_en.tex  (run twice for references)
\documentclass[conference]{IEEEtran}

\usepackage{amssymb}
\usepackage{amsmath}
\usepackage{booktabs}
\usepackage{graphicx}
\usepackage{url}
\usepackage{microtype}

\begin{document}

\title{DeReFusion Ablation Study Report: Component Contributions and a\\
Weight-Initialization Confound}

\author{
\IEEEauthorblockN{Geng}
\IEEEauthorblockA{Applied Soft Computing study project\\
Email: gengyifu686@users.noreply.github.com}
}

\maketitle

\begin{abstract}
This report documents a controlled component ablation of the DeReFusion
forecasting framework. Three ablation variants---\texttt{woDy} (base branch
only), \texttt{woLSTM} (Transformer-only residual), and
\texttt{woTransformer} (LSTM-only residual)---are evaluated on two assets
(S\&P~500 index and BTC/USD) at horizon $T{=}24$, seed 2021, under the
protocol of the accompanying reproduction study. Seven runs are reported:
the three variants on each of two assets, plus a \texttt{revin-DLinear}
environment-calibration run whose result (MSE 0.07007) matches the reference
value exactly. Three findings stand out. First, the LSTM is the dominant
component of the residual branch: removing it \emph{improves} the model on
GSPC (0.06150 vs.\ 0.06230) and the Transformer-only variant is the best of
the three on both assets. Second, the \texttt{woDy} ablation is confounded:
it differs from the \texttt{revin-DLinear} baseline not by removing a
component but by dropping the DLinear weight-initialization convention, and
this alone accounts for an 8\% MSE gap at identical parameter count. Third,
component value is asset-dependent: the full hybrid wins on BTC/USD while
the LSTM-free variant wins on GSPC. All seven results are reproduced from
the official code on CPU with a single seed.
\end{abstract}

\begin{IEEEkeywords}
financial time series forecasting, ablation study, DeReFusion,
decomposition--residual architecture, weight initialization,
reversible instance normalization
\end{IEEEkeywords}

\section{Introduction}

DeReFusion combines a linear decomposition base branch (DLinear) with a
learned LSTM--Transformer residual branch under direct additive fusion,
wrapped in Reversible Instance Normalization (RevIN). The accompanying
reproduction study establishes the pipeline and its central qualitative
claim; this report isolates the question that follows immediately: which
components of the residual branch actually carry the model, and what does
each removal cost?

The study is deliberately narrow. It covers two assets (GSPC, BTCUSD), one
horizon ($T{=}24$), one seed (2021), and three ablation variants, for seven
executed runs in total. Within that scope the results are unambiguous and,
in one case, counter-intuitive: the most complex ablation is not the best
one, and one of the ablation variants is not a clean ablation at all.

The remainder is organized as follows. Section~\ref{sec:design} defines the
variants and the controlled protocol. Section~\ref{sec:setup} records the
environment and the calibration run. Section~\ref{sec:results} presents the
measured results. Section~\ref{sec:analysis} interprets them, including the
initialization confound. Section~\ref{sec:conclusion} concludes and lists
next steps.

\section{Ablation Design}
\label{sec:design}

\subsection{Variants}

The repository ships three ablation models under
\texttt{models/derefusion/ablation\_variant/}. All three retain RevIN, the
DLinear base branch, and direct additive fusion; they differ only in the
residual branch or in its removal:

\begin{enumerate}
  \item \textbf{\texttt{woDy} --- base only.} The residual branch is removed
        entirely; the model reduces to RevIN-wrapped DLinear
        (4{,}664 parameters).
  \item \textbf{\texttt{woLSTM} --- Transformer residual.} The residual
        branch keeps a single Transformer encoder layer (4 heads, FFN
        $4d$, dropout 0.3) and drops the LSTM (19{,}988 parameters).
  \item \textbf{\texttt{woTransformer} --- LSTM residual.} The residual
        branch keeps the LSTM and drops the Transformer
        (11{,}988 parameters).
\end{enumerate}

For reference, the full DeReFusion residual branch (LSTM $\rightarrow$
Transformer) has ${\sim}28$k parameters.

\subsection{Controlled Protocol}

Every run uses an identical protocol, unchanged from the reproduction
baseline: chronological 70/10/20 split; StandardScaler fitted on the
training split only; RevIN on every trained model; multivariate input with
single-target output (MS; target \texttt{Close}); look-back $L{=}96$;
label length 48; Adam, learning rate $10^{-4}$, batch size 32, MSE loss;
cosine schedule; early stopping with patience 5 on validation MSE;
${\sim}30$ epochs; seed 2021; CPU execution. Metrics are MSE, MAE, RMSE,
MAPE, MSPE, and $R^2$, computed on the normalized series.

The grid is $3$ variants $\times$ $2$ assets $=6$ runs, plus one
calibration run (Section~\ref{sec:setup}), for seven runs total.

\section{Environment and Calibration}
\label{sec:setup}

All runs execute on a single Windows machine with a CPU-only PyTorch build,
no WSL, no GPU. Table~\ref{tab:config} lists the shared configuration.

\begin{table}[!t]
\centering
\caption{Shared configuration of all ablation runs.}
\label{tab:config}
\small
\begin{tabular}{@{}p{2.9cm}p{4.0cm}@{}}
\toprule
Parameter & Value \\
\midrule
Look-back $L$ / label len. & 96 / 48 \\
Horizon $T$ & 24 \\
Input channels / target & 4 (OHLC) / Close \\
$d_{\mathrm{model}}$ & 32 \\
Moving-average kernel $k$ & 25 \\
Residual Transformer / LSTM & 1 layer each (4 heads, FFN $4d$) \\
Dropout & 0.3 \\
Optimizer / learning rate & Adam / $1{\times}10^{-4}$ \\
Batch size / loss & 32 / MSE \\
Schedule & cosine \\
Early stopping & patience 5 (val.\ MSE) \\
Random seed & 2021 \\
Device & CPU \\
\bottomrule
\end{tabular}
\end{table}

\paragraph{Calibration run.} Following the receipt protocol, a
\texttt{revin-DLinear} baseline with a known reference answer was executed
first: GSPC, $T{=}24$, seed 2021. The measured MSE is
\textbf{0.070065}, against a reference of 0.07007 --- a difference of
$6\times10^{-6}$, far inside the accepted tolerance band of 0.066--0.074.
The environment is therefore confirmed correct, and the ablation results
below can be compared directly with the reference table.

\section{Results}
\label{sec:results}

\subsection{Measured Metrics}

Table~\ref{tab:main} reports the seven executed runs. All values are
computed on the normalized series, consistent with the paper's convention;
the best MSE within each asset group is bold.

\begin{table*}[!t]
\centering
\caption{Ablation results, all runs executed locally ($T{=}24$, seed 2021,
CPU). Best MSE per asset group in bold. MAPE/MSPE in percent.}
\label{tab:main}
\small
\begin{tabular}{llccccccc}
\toprule
Asset & Model & MSE & MAE & RMSE & MAPE (\%) & MSPE (\%) & $R^2$ & Params \\
\midrule
GSPC   & \texttt{woLSTM}        & \textbf{0.06150} & 0.18714 & 0.24799 & 5.80 & 0.651 & 0.8708 & 19{,}988 \\
GSPC   & \texttt{woTransformer} & 0.06897 & 0.20742 & 0.26262 & 6.36 & 0.704 & 0.8551 & 11{,}988 \\
GSPC   & \texttt{revin-DLinear} & 0.07007 & 0.21356 & 0.26470 & 6.52 & 0.712 & 0.8528 & 4{,}664 \\
GSPC   & \texttt{woDy}          & 0.07586 & 0.21988 & 0.27543 & 6.71 & 0.765 & 0.8406 & 4{,}664 \\
\midrule
BTCUSD & \texttt{woLSTM}        & \textbf{0.22178} & 0.36189 & 0.47094 & 9.32 & 1.504 & 0.8734 & 19{,}988 \\
BTCUSD & \texttt{woTransformer} & 0.22942 & 0.36869 & 0.47898 & 9.45 & 1.539 & 0.8690 & 11{,}988 \\
BTCUSD & \texttt{woDy}          & 0.23727 & 0.37389 & 0.48711 & 9.58 & 1.583 & 0.8645 & 4{,}664 \\
\bottomrule
\end{tabular}
\end{table*}

\subsection{Comparison With the Full Model}

Table~\ref{tab:ref} places the best ablation per asset against the full
DeReFusion model and the linear baseline. The reference values for the full
model and the baseline are reported by the collaborating study on the same
data and protocol; they are reproduced here for comparison and are not
re-executed in this work.

\begin{table}[!t]
\centering
\caption{Best ablation vs.\ reference full model and linear baseline
($T{=}24$, seed 2021). Lower MSE is better.}
\label{tab:ref}
\small
\begin{tabular}{@{}p{1.5cm}p{2.9cm}cc@{}}
\toprule
Asset & Model & MSE & $R^2$ \\
\midrule
GSPC   & \textbf{woLSTM} (local)   & \textbf{0.06150} & 0.8708 \\
GSPC   & DeReFusion (reference)    & 0.06230 & --- \\
GSPC   & revin-DLinear (local)     & 0.07007 & 0.8528 \\
\midrule
BTCUSD & DeReFusion (reference)    & 0.21797 & --- \\
BTCUSD & \textbf{woLSTM} (local)   & \textbf{0.22178} & 0.8734 \\
BTCUSD & revin-DLinear (ref.)      & 0.22322 & --- \\
\bottomrule
\end{tabular}
\end{table}

\section{Analysis and Findings}
\label{sec:analysis}

\subsection{Finding 1: The LSTM Is Not the Component That Carries GSPC}

The clearest result in Table~\ref{tab:main} is that removing the LSTM
\emph{improves} GSPC: \texttt{woLSTM} reaches MSE 0.06150, better than the
full model's 0.06230 and better than every other variant. A three-point
ordering emerges on GSPC:

\begin{equation}
\texttt{woLSTM} \;<\; \texttt{woTransformer} \;<\; \texttt{revin-DLinear} \;\approx\; \texttt{woDy}.
\end{equation}

Interpreted component-wise, the Transformer contributes positively to the
residual branch (its removal costs 0.0075 MSE relative to \texttt{woLSTM}),
while the LSTM contributes negatively --- adding it back raises MSE slightly.
On a near-random-walk equity index, the LSTM's local sequential memory
appears to fit noise rather than signal at this sample size.

This is a negative result and is reported as such: the ablation does not
show that every component helps. On BTCUSD the picture reverses
(Section~\ref{sec:analysis-C}), so the negative LSTM contribution is
asset-specific rather than universal.

\subsection{Finding 2: The \texttt{woDy} Ablation Is Confounded by Initialization}
\label{sec:analysis-B}

\texttt{woDy} is intended as the pure base-branch ablation, and by
construction its forward computation is identical to the
\texttt{revin-DLinear} baseline: RevIN normalize, moving-average
decomposition, two linear maps, add, RevIN de-normalize. Both models have
exactly 4{,}664 parameters. Yet on GSPC they differ by 8\% in MSE
(0.07586 vs.\ 0.07007).

Inspection of the two implementations locates the cause. The
\texttt{revin-DLinear} model follows the DLinear paper's convention and
initializes both linear layers to the uniform averaging kernel:

\begin{equation}
W_{\text{se}}, W_{\text{tr}} \;\leftarrow\; \tfrac{1}{L}\,\mathbf{1}_{T\times L},
\end{equation}

which makes the model start as a naive moving-average extrapolator.
\texttt{woDy} uses PyTorch's default Kaiming-uniform initialization for the
same layers. Same architecture, same parameter count, different starting
point --- and a materially different optimum.

The practical consequence is that \texttt{woDy} is \emph{not} a clean
measure of the residual branch's contribution: its 0.07586 gap to the full
model mixes the removal of the residual branch with the loss of the
initialization convention. A clean base-only ablation must reuse the
DLinear initialization; we flag this as a correctness issue rather than
report it as a component finding.

\subsection{Finding 3: Component Value Is Asset-Dependent}
\label{sec:analysis-C}

On BTCUSD the ordering is different. \texttt{woLSTM} still leads the
ablations, but now the full model wins: DeReFusion at 0.21797 beats
\texttt{woLSTM} at 0.22178, which in turn beats the baseline at 0.22322.
The LSTM, which hurt on GSPC, helps on BTCUSD --- consistent with an asset
whose dynamics are less random-walk-like and therefore carry learnable
sequential structure.

Across both assets the Transformer-only variant is more stable than the
LSTM-only variant, and the LSTM-only residual is consistently the weakest
of the two residual designs where both are present. The residual branch's
value is thus real but asset-conditioned, which is exactly the regime that
motivates structure-aware routing.

\subsection{Parameter Efficiency}

Parameter count does not track accuracy. \texttt{woLSTM} (19{,}988
parameters) beats \texttt{woTransformer} (11{,}988) and \texttt{woDy}
(4{,}664) on both assets, but the ranking is not monotone in size:
\texttt{woDy} matches the baseline's parameter count while performing
worse, because of the confound above. Within the residual designs, the
larger variant wins on both assets.

\section{Conclusion and Next Steps}
\label{sec:conclusion}

We executed a seven-run the component ablation of DeReFusion on GSPC and
BTCUSD at $T{=}24$, seed 2021. Three conclusions follow.

\begin{enumerate}
  \item The LSTM is not universally load-bearing. Removing it improves
        GSPC and degrades BTCUSD, so the residual branch's value is
        asset-dependent.
  \item The \texttt{woDy} variant is confounded by weight initialization
        and must not be read as a clean measure of the residual branch's
        contribution. The fix is to restore the DLinear
        $1/L$ initialization in the ablation.
  \item The Transformer-only residual is the best ablation on both assets,
        suggesting it is the more robust residual design at this data
        scale.
\end{enumerate}

\paragraph{Next steps.} Re-run \texttt{woDy} with the DLinear
initialization to obtain an unconfounded base-only number; extend the grid
to the remaining horizons ($T\in\{1,7,12,36\}$) and to the additional
assets; and add a second seed to separate the two very close GSPC
comparisons (\texttt{woLSTM} vs.\ full model, a 1.3\% gap) from seed noise.

\section*{Acknowledgment}

The ablation models, the dataset, and the training harness come from the
DeReFusion repository, released under an MIT license.

\begin{thebibliography}{9}

\bibitem{derefusion}
C.-C. Hsieh and M.-Y. Chen,
``DeReFusion: A controlled comparison of soft computing fusion strategies
for financial time series forecasting via a decomposition-residual
architecture,'' \emph{Applied Soft Computing}, vol. 203, p. 116252, 2026,
doi: 10.1016/j.asoc.2026.116252.

\bibitem{kim2022revin}
T. Kim \emph{et al.}, ``Reversible instance normalization for accurate
time-series forecasting against distribution shift,'' in \emph{Proc. ICLR},
2022.

\bibitem{zeng2023dlinear}
A. Zeng, M. Chen, L. Zhang, and Q. Xu, ``Are transformers effective for time
series forecasting?'' in \emph{Proc. AAAI}, 2023.

\end{thebibliography}

\end{document}
